\section{What is verified, and how}
\label{sec:verification}

\subsection{No gaps}
\label{sec:nogaps}

A search of all 2,659 Lean files for the token \lean{sorry} returns five occurrences.
One is inside a docstring. The other four are in
\file{ComparatorChallenges/NavierStokes.lean} (lines 277 and 284) and
\file{ComparatorChallenges/Euler.lean} (lines 88 and 184), i.e.\ in the \emph{challenge
statements} that reproduce the Formal Conjectures placeholders for the Millennium problem.
Those four are not used by any proof: the challenge modules import only \Mathlib, and
nothing imports them. The two proof libraries contain no \lean{sorry}. A search for
\lean{axiom}, \lean{native\_decide}, \lean{implemented\_by}, \lean{@[extern]},
\lean{unsafe}, \lean{partial} and \lean{opaque} finds none in \file{NavierStokes/} or
\file{Euler/} either; the occurrences of the word ``axiom'' in those directories are
docstrings explaining that something is \emph{not} an axiom, for example
\file{NavierStokes/ProblemStatement.lean}: ``There is no proof, witness, or axiom
asserting this proposition in this module.''

\subsection{Axioms}
\label{sec:axioms}

The four headline declarations carry explicit \lean{\#print axioms} commands in the
repository: \file{NavierStokes/ComparatorSolution.lean} lines 31--32 for the two
Navier--Stokes alternatives, and \file{Euler/Solution.lean} lines 73 and 75 for the two
Euler statements. The reported dependency sets are
\[
  \{\lean{propext},\ \lean{Classical.choice},\ \lean{Quot.sound}\},
\]
the standard classical base of \Mathlib{} analysis. Section~\ref{sec:migration-results}
reports the same commands re-run after the migration.

We re-ran those commands at the pinned toolchain after the migration, and can now report
the compiler's own answer rather than the repository's claim:

\begin{lstlisting}[language=ML]
'NavierStokesR3.breakdownStatement' depends on axioms: [propext, Classical.choice, Quot.sound]
'NavierStokes.PeriodicPaper.breakdownStatement' depends on axioms: [propext, Classical.choice, Quot.sound]
'NavierStokes.Comparator.navier_stokes_breakdown_R3' depends on axioms: [propext, Classical.choice, Quot.sound]
'NavierStokes.Comparator.navier_stokes_breakdown_periodic' depends on axioms: [propext, Classical.choice, Quot.sound]
\end{lstlisting}

This matters more than it might appear. A formalization of this size has many places
where a proof could in principle be closed by a stronger assumption---an inverse function
theorem asserted as an axiom, a well-posedness theorem assumed for the class at hand. The
axiom audit is what rules that out: an unproved input would show up as an extra axiom in
the dependency set of the headline theorem, and there is none.

\subsection{Statement fidelity}
\label{sec:fidelity}

Typechecking a theorem says nothing about whether it states the intended mathematics. We
checked the statements in three ways.

\emph{Definitions copied, not paraphrased.} The Euler definitions
(\file{Euler/SolutionDefinitions.lean}) and both challenge modules are adapted from the
Formal Conjectures formalization of the Millennium problem statement, with a comment
recording the upstream commit. The task was therefore to prove \emph{that} statement, not
a convenient variant.

\emph{Every hypothesis is explicit.} The Navier--Stokes statement's force class, for
example, is \lean{CompactPositiveTimeSupport f}, defined as
\lean{HasCompactSupport f ∧ tsupport f ⊆ Ioi 0 ×ˢ univ}: smooth, compactly supported in
spacetime, supported at strictly positive times. The energy condition is
\lean{UniformFiniteEnergy (Ico 0 1) u}, which asks for a single $E\ge0$ with
\lean{SquareIntegrableAtTime u t} and \lean{kineticEnergy u t ≤ E} for every $t\in[0,1)$.
Neither condition can be satisfied vacuously, and neither hides an assumption about the
\emph{competitor}, which is quantified with no support, growth or energy-inequality
hypothesis beyond the bound.

\emph{No statement is vacuous.} The two constructions are explicit: the repository builds
the fields and then proves the properties, so the existential statements have witnesses
that can be inspected. Theorems that are only \emph{conditional} are named as such; for
example \file{NavierStokes/ProblemStatement.lean} contains the proposition
\lean{candidateStatement} together with a docstring saying that nothing in the module
proves it, and the proof that it holds is in \file{NavierStokes/R3/Theorem.lean} under the
same name. This is a good convention: a reader who searches for a statement finds, at the
statement, a pointer to where it is discharged.

\subsection{Independent checking}
\label{sec:comparator}

The repository ships a Comparator configuration (\file{ComparatorChallenges/}), which
exports the four headline theorems to Lean's kernel-independent proof format and replays
them through a separate checker after translation. This is a defence against a bug in
Lean's own kernel or in the elaborator, and it is the strongest available form of
independent verification for a Lean development. The migration described in
Section~\ref{sec:libraries} forced a choice here: Comparator is released per toolchain and
its releases stop at \Lean{} \texttt{v4.35.0}, one version below the toolchain that
\TauCeti{} requires. We kept the migration and the challenge definitions and dropped the
Comparator \emph{library} dependency; the challenge files remain in the repository and can
be re-checked with a Comparator installation matching whichever toolchain is in use.
